\documentclass[../main.tex]{subfiles}
\ifSubfilesClassLoaded{\setcounter{chapter}{12}}{}
\begin{document}

\chapter{Preprocessor, \texttt{Makefile}, dan \texttt{gdb} Pengenalan}

\begin{subcpmk}
  \item Sub-CPMK 7.4: Preprocessor directive: \texttt{\#define}, \texttt{\#include}, makro dengan parameter, kompilasi kondisional, include guards, predefined macros
  \item Sub-CPMK 6.2 (lanjutan): Target \texttt{Makefile} dan otomasi build
  \item Sub-CPMK 6.3 (lanjutan): \texttt{gdb}: breakpoint, \texttt{run}, \texttt{backtrace} pada program kecil
\end{subcpmk}

\SesiRPS{13}{7.4, 6.2--6.3}{$\sim$170 menit}{CPMK-7}

% ============================================================
\section{Sarana dan prasarana}
% ============================================================
\texttt{make} (atau \texttt{mingw32-make} di Windows), \texttt{gdb}, dokumentasi GCC.

% ============================================================
\section{Preprocessor: Apa dan Mengapa}
% ============================================================

Preprocessor adalah tahap \textbf{pertama} dalam proses kompilasi C --- ia memproses semua baris yang diawali tanda \texttt{\#} \textbf{sebelum} kompilator melihat kode. Preprocessor melakukan manipulasi teks: menyertakan isi file header, mengganti makro dengan nilainya, dan memasukkan atau mengeluarkan blok kode berdasarkan kondisi. Outputnya adalah kode C ``murni'' yang kemudian dikompilasi.

Memahami preprocessor penting karena banyak konstruksi C bergantung padanya --- termasuk \texttt{\#include} yang kita gunakan di setiap program, \texttt{\#define} untuk konstanta, include guards untuk mencegah penyertaan ganda, dan kompilasi kondisional untuk membuat kode portabel antar platform.

% ============================================================
\section{Tabel Directive Preprocessor}
% ============================================================

{\small
\begin{tabular}{|>{\raggedright\arraybackslash}p{3.6cm}|>{\raggedright\arraybackslash}p{9.4cm}|}
\hline
\textbf{Directive} & \textbf{Keterangan} \\
\hline
\texttt{\#include <file>} & Menyertakan header sistem (dicari di path sistem) \\
\hline
\texttt{\#include "file"} & Menyertakan header lokal (dicari di direktori saat ini dulu) \\
\hline
\texttt{\#define NAMA nilai} & Mendefinisikan makro (penggantian teks) \\
\hline
\texttt{\#define NAMA(x) ((x)*(x))} & Makro dengan parameter (function-like macro) \\
\hline
\texttt{\#undef NAMA} & Menghapus definisi makro \\
\hline
\texttt{\#ifdef NAMA} / \texttt{\#ifndef} & Kompilasi kondisional berdasarkan ada/tidaknya makro \\
\hline
\texttt{\#if expr} / \texttt{\#elif} / \texttt{\#else} / \texttt{\#endif} & Kompilasi kondisional berdasarkan ekspresi \\
\hline
\texttt{\#pragma once} & Include guard non-standar (didukung mayoritas compiler) \\
\hline
\texttt{\#error "pesan"} & Menghasilkan error kompilasi dengan pesan kustom \\
\hline
\end{tabular}
}

% ============================================================
\section{Makro \texttt{\#define}}
% ============================================================

\subsection{Makro Objek (Konstanta)}

\begin{lstlisting}[language=C,caption=Makro konstanta]
#define MAX_SISWA   40
#define PI          3.14159265
#define NAMA_PRODI  "Teknik Informatika"

printf("Prodi: %s, Max siswa: %d\n", NAMA_PRODI, MAX_SISWA);
\end{lstlisting}

\subsection{Makro Fungsi (Dengan Parameter)}

Makro fungsi melakukan penggantian teks dengan parameter. \textbf{Selalu bungkus} parameter dan seluruh ekspresi dalam tanda kurung untuk menghindari bug dari precedence operator.

\begin{lstlisting}[language=C,caption=Makro fungsi --- benar dan berbahaya]
// BENAR: parameter dan ekspresi dibungkus kurung
#define KUADRAT(x) ((x) * (x))
#define MAKS(a, b) ((a) > (b) ? (a) : (b))

// BERBAHAYA: tanpa kurung
#define BURUK(x) x * x
// BURUK(1+2) -> 1+2 * 1+2 = 1+2+2 = 5 (bukan 9!)

#include <stdio.h>
int main(void) {
    printf("KUADRAT(1+2) = %d\n", KUADRAT(1+2));  // 9 (benar)
    printf("BURUK(1+2) = %d\n", BURUK(1+2));       // 5 (salah!)
    printf("MAKS(3, 7) = %d\n", MAKS(3, 7));       // 7
    return 0;
}
\end{lstlisting}

\begin{catatan}
\textbf{Makro vs Fungsi:} Makro lebih cepat (tidak ada overhead pemanggilan) tetapi tidak memiliki pengecekan tipe, bisa menyebabkan efek samping (\texttt{KUADRAT(x++)} mengevaluasi \texttt{x++} dua kali!), dan lebih sulit di-debug. Untuk kode modern, lebih disarankan menggunakan fungsi \texttt{inline} (C99) atau fungsi biasa.
\end{catatan}

% ============================================================
\section{Kompilasi Kondisional}
% ============================================================

Kompilasi kondisional memungkinkan blok kode dimasukkan atau dikeluarkan dari kompilasi berdasarkan kondisi tertentu. Ini sangat berguna untuk: kode debug yang hanya aktif saat pengembangan, kode portabel antar platform (Windows vs Linux), dan fitur opsional.

\begin{lstlisting}[language=C,caption=Kompilasi kondisional]
#include <stdio.h>

#define DEBUG  // aktifkan mode debug (hapus untuk menonaktifkan)

int main(void) {
    int x = 42;
    
    #ifdef DEBUG
        printf("[DEBUG] x = %d, alamat = %p\n", x, (void*)&x);
    #endif
    
    printf("Nilai: %d\n", x);
    
    #if defined(_WIN32)
        printf("Platform: Windows\n");
    #elif defined(__linux__)
        printf("Platform: Linux\n");
    #elif defined(__APPLE__)
        printf("Platform: macOS\n");
    #else
        printf("Platform: tidak dikenal\n");
    #endif
    
    return 0;
}
\end{lstlisting}

% ============================================================
\section{Predefined Macros}
% ============================================================

\begin{tabular}{|l|p{8cm}|}
\hline
\textbf{Makro} & \textbf{Isi} \\
\hline
\texttt{\_\_FILE\_\_} & Nama file sumber saat ini (string) \\
\hline
\texttt{\_\_LINE\_\_} & Nomor baris saat ini (integer) \\
\hline
\texttt{\_\_DATE\_\_} & Tanggal kompilasi (string: ``Apr 20 2026'') \\
\hline
\texttt{\_\_TIME\_\_} & Waktu kompilasi (string: ``14:30:00'') \\
\hline
\texttt{\_\_func\_\_} & Nama fungsi saat ini (C99+) \\
\hline
\end{tabular}

\begin{lstlisting}[language=C,caption=Menggunakan predefined macros untuk log]
#include <stdio.h>

#define LOG(msg) printf("[%s:%d %s] %s\n", \
    __FILE__, __LINE__, __func__, msg)

int main(void) {
    LOG("Program dimulai");
    
    int x = 10;
    LOG("Variabel x diinisialisasi");
    
    printf("Dikompilasi: %s %s\n", __DATE__, __TIME__);
    return 0;
}
\end{lstlisting}

% ============================================================
\section{Makefile: Otomasi Build}
% ============================================================

\texttt{Makefile} adalah skrip untuk program \texttt{make} yang mengotomasi kompilasi --- menentukan file mana yang bergantung pada file lain dan perintah apa yang dijalankan. Pada proyek multi-file, \texttt{Makefile} memastikan hanya file yang berubah yang dikompilasi ulang, menghemat waktu secara signifikan.

\begin{lstlisting}[language=make,caption=Makefile dengan penjelasan]
# Variabel
CC = gcc
CFLAGS = -Wall -Wextra -std=c17 -g
TARGET = program
SRCS = main.c mathutil.c
OBJS = main.o mathutil.o    # .c -> .o

# Target default
all: $(TARGET)

# Linking: gabungkan semua .o
$(TARGET): $(OBJS)
	$(CC) $(CFLAGS) -o $@ $^

# Kompilasi: .c -> .o (aturan implisit)
%.o: %.c
	$(CC) $(CFLAGS) -c $< -o $@

# Dependency header
main.o: main.c mathutil.h
mathutil.o: mathutil.c mathutil.h

# Bersihkan
clean:
	rm -f $(TARGET) $(OBJS)

# Jalankan program
run: $(TARGET)
	./$(TARGET)

.PHONY: all clean run
\end{lstlisting}

\begin{tabular}{|l|l|}
\hline
\textbf{Simbol Make} & \textbf{Arti} \\
\hline
\texttt{\$@} & Nama target \\
\hline
\texttt{\$<} & Prerequsite pertama \\
\hline
\texttt{\$\^{}} & Semua prerequsite \\
\hline
\texttt{\%.o: \%.c} & Pattern rule: setiap \texttt{.o} dari \texttt{.c} yang sesuai \\
\hline
\texttt{.PHONY} & Target yang bukan file (selalu dieksekusi) \\
\hline
\end{tabular}

\begin{catatan}
Pada Windows dengan MinGW, gunakan \texttt{mingw32-make} sebagai pengganti \texttt{make}. Target \texttt{clean} mungkin perlu menggunakan \texttt{del} alih-alih \texttt{rm -f}.
\end{catatan}

% ============================================================
\section{GDB: Debugger Pengenalan}
% ============================================================

\texttt{gdb} (\textit{GNU Debugger}) memungkinkan kita menjalankan program langkah demi langkah, memeriksa nilai variabel, dan menemukan sumber bug. Untuk menggunakan \texttt{gdb}, kompilasi program dengan opsi \texttt{-g} (menyertakan informasi debug).

\begin{tabular}{|l|p{8cm}|}
\hline
\textbf{Perintah GDB} & \textbf{Keterangan} \\
\hline
\texttt{gdb ./program} & Mulai GDB dengan program yang akan di-debug \\
\hline
\texttt{break main} & Set breakpoint di awal fungsi \texttt{main} \\
\hline
\texttt{break file.c:15} & Set breakpoint di baris 15 \\
\hline
\texttt{run} / \texttt{r} & Jalankan program \\
\hline
\texttt{next} / \texttt{n} & Jalankan satu baris (lewati fungsi) \\
\hline
\texttt{step} / \texttt{s} & Jalankan satu baris (masuk ke fungsi) \\
\hline
\texttt{print x} / \texttt{p x} & Tampilkan nilai variabel \texttt{x} \\
\hline
\texttt{backtrace} / \texttt{bt} & Tampilkan call stack \\
\hline
\texttt{continue} / \texttt{c} & Lanjutkan sampai breakpoint berikutnya \\
\hline
\texttt{quit} / \texttt{q} & Keluar dari GDB \\
\hline
\texttt{list} / \texttt{l} & Tampilkan kode sumber di sekitar posisi saat ini \\
\hline
\texttt{info locals} & Tampilkan semua variabel lokal \\
\hline
\texttt{watch x} & Berhenti ketika nilai \texttt{x} berubah \\
\hline
\end{tabular}

\begin{contoh}
\textbf{Contoh sesi GDB:}

\begin{verbatim}
$ gcc -Wall -Wextra -std=c17 -g -o program main.c
$ gdb ./program
(gdb) break main
Breakpoint 1 at 0x...: file main.c, line 5.
(gdb) run
Starting program: ./program
Breakpoint 1, main () at main.c:5
5       int x = 10;
(gdb) next
6       int y = 20;
(gdb) print x
$1 = 10
(gdb) next
7       int z = x + y;
(gdb) print y
$2 = 20
(gdb) next
(gdb) print z
$3 = 30
(gdb) backtrace
#0  main () at main.c:8
(gdb) quit
\end{verbatim}
\end{contoh}

% ============================================================
\section{Referensi sesi}
% ============================================================
\begin{itemize}[leftmargin=*]
  \item GCC manual, preprocessor. \url{https://gcc.gnu.org/onlinedocs/gcc/index.html}
  \item GNU C Reference, preprocessor. \url{https://www.gnu.org/software/gnu-c-manual/gnu-c-manual.html}
  \item GNU Make Manual. \url{https://www.gnu.org/software/make/manual/make.html}
  \item GDB Documentation. \url{https://sourceware.org/gdb/current/onlinedocs/gdb/}
  \item UNWAHAS JTI2403 (struktur OBE praktikum). \url{https://simpel.teknik.unwahas.ac.id/rps/191}
\end{itemize}

% ============================================================
\begin{aktivitas}
  \item Buat \texttt{Makefile} dengan target \texttt{all}, \texttt{clean}, dan \texttt{run} untuk proyek multi-file.
  \item Kompilasi dengan \texttt{-g} dan lakukan sesi GDB: set breakpoint, jalankan langkah demi langkah, periksa variabel.
  \item Buat makro \texttt{LOG} yang mencetak file, baris, dan pesan.
  \item Buat program dengan \texttt{\#ifdef DEBUG} untuk mencetak informasi debug opsional.
\end{aktivitas}

\begin{latihan}
  \item Mengapa makro \texttt{\#define KUADRAT(x) x*x} berbahaya? Berikan contoh.
  \item Apa keuntungan kompilasi terpisah (\texttt{.c} $\rightarrow$ \texttt{.o} $\rightarrow$ executable)?
  \item Jelaskan perbedaan \texttt{next} dan \texttt{step} di GDB.
  \item Bagaimana cara membuat Makefile yang hanya mengompilasi ulang file yang berubah?
\end{latihan}

\begin{asesmen}
\textbf{Tugas M11:} proyek kecil multi-berkas (min.\ 3 file) + \texttt{Makefile} fungsional + cuplikan sesi \texttt{gdb} (perintah dan keluaran, screenshot terminal).
\end{asesmen}

\begin{checklist}
  \item Saya memahami directive preprocessor utama
  \item Saya dapat menulis makro yang aman (dengan kurung)
  \item Saya dapat membuat \texttt{Makefile} fungsional
  \item Saya dapat menjalankan sesi debug dasar dengan \texttt{gdb}
  \item Saya memahami kompilasi kondisional (\texttt{\#ifdef}/\texttt{\#endif})
\end{checklist}

\begin{rangkuman}
Bab ini membahas preprocessor (directive utama, makro objek dan fungsi, kompilasi kondisional, predefined macros), otomasi build dengan \texttt{Makefile} (variabel, target, pattern rules), dan pengenalan debugger \texttt{gdb} (breakpoint, step, print, backtrace).

\textbf{Poin kunci:} selalu kurung parameter makro; \texttt{Makefile} menghemat waktu kompilasi; \texttt{-g} wajib untuk debug; \texttt{gdb} membantu menemukan bug tanpa \texttt{printf} debugging.

\textbf{Kata kunci:} preprocessor, \texttt{\#define}, \texttt{\#ifdef}, makro, \texttt{Makefile}, \texttt{make}, \texttt{gdb}, breakpoint, debugging
\end{rangkuman}

\end{document}
